\cotefile{157-5}
\scanpage{341}
% top centre
\mspage{257}
% lines 1–2
Ainsi, pour définir \inserted{les} $X_{\alpha+1}(j)$, il faut \deleted{définir} tenir compte de toutes ces contraintes.

% line 3
Donc

% lines 4–5
$1^\circ)$ Pour $j^{\unsure{\prime\prime|\sharp}}$ fixé, $X_{\alpha+1}(j'')$ doit \deleted{\unsure{recevoir}}\inserted{être but de} des flèches
% line 6
\[
 X_\alpha(j') \unsure{\xrightarrow{t_\alpha(\lambda)}|\overset{t_\alpha(\lambda)}{\rule{2.8em}{0.4pt}}} X_{\alpha+1}(j''),
\]
% lines 7–10
pour les $\lambda:j'\unsure{\longrightarrow|\mathrel{-}}j''$ dans $\operatorname{Ob}J_{/j''}$ -- y compris pour $\lambda=\mathrm{id}_{j''}$, auquel cas $t_\alpha(\lambda)$\inserted{$:X_\alpha(j'')\longrightarrow X_{\alpha+1}(j'')$} \unsure{jouera|joue} le rôle du morphisme de transition $\tau_\alpha(j'')$.

% lines 11–18
$\uline{2}^\circ)$ Si $\alpha$ \deleted{n'est pas \unsure{ord$^{\mathrm{al}}$} limite, i.e. de} \inserted{est de la forme $\alpha'+2$, \deleted{alors}} \deleted{\unsure{ou|soit}} \inserted{i.e} s'il n'est \deleted{pas}\inserted{ni} \unsure{un ord$^{\mathrm{inal}}$} limite, ni \deleted{\unsure{de} la forme $\unsure{\alpha'+1}$, alors les flèches \unsure{qui|prit} la forme $\unsure{\beta+1}$ on peut mettre $\tau_\alpha(\unsure{j'})$} successeur d'un ordinal limite, alors on \unsure{exige} de plus, pour les morphismes de transition $\tau_\alpha(j\unsure{'|''})$, \deleted{\unsure{resp.} $\tau_\alpha(j'')$} d'égaliser les \inserted{couples de} flèches \unsure{prévues|prises} dans \unsure{ii)|(ii)} et (iii) ci-dessus.

% lines 19–20
Comme \unsure{l'ens|l'une} de ces données du $1^\circ)$, ou de ces contraintes du $\uline{2}^\circ)$,
% archive sheet number: 340
